% !TeX root = ../nc-quadric-surface-moduli.tex
% 2026-11-21 / v1.34（論文Bへ証明移管）
\section{幾何的二周期 cubic 族の普遍 canonical hull}
\label{sec:universal-cubic-hull}
\paragraph{主張の範囲.}
任意の幾何的二周期 cubic 族について、全幾何ファイバーの canonical hull を
有限 projective 加群間の射で相対化する。
可換二次曲面の特殊性、Artin 厚み、Krull 交叉定理を用いない。
体上の存在定理は Van den Bergh の Theorem~6.1 に依存し、
本節ではそこで存在する hull の六次元関係空間を、逆 coevaluation の核として
正確に同定する。相対化の新規性は未確認である。
二周期の存在と、その二乗が反 Serre 関手であることは別条件である。

\subsection{入力と定理}

$k$ を標数零の体、$R$ を可換 Noetherian $k$-代数とする。
$D$ を正の $R$-線形 $\mathbb Z$-代数とし、次を仮定する。
\begin{enumerate}
\item $D_{ii}=R$、$D_{ij}=0$ $(j<i)$、全ての $D_{ij}$ は有限 projective である。
  $V_i=D_{i,i+1}$ の階数は二であり、$D$ は三次関係の直和因子
  $R_i\subset V_i\otimes V_{i+1}\otimes V_{i+2}$、$\rank R_i=2$ により提示される。
\item 四次の交わり
  \[
  W_i=(R_i\otimes V_{i+3})\cap(V_i\otimes R_{i+1})
  \]
  は可逆 $R$-加群である。交わりは任意の基底変更と可換し、縮約により
  \[
  W_i\otimes V_{i+3}^{\vee}\xrightarrow{\sim}R_i,
  \qquad W_i\otimes V_i^{\vee}\xrightarrow{\sim}R_{i+1}
  \tag{UH-1}\label{hull-eq:coevaluation}
  \]
  が得られる。
\item 全幾何ファイバーは geometric quintuple に対応する標準 cubic AS 正則
  $\mathbb Z$-代数である。標準 cubic 解消は各内部次数で有限 projective の
  分裂完全複体となり、任意の基底変更に整合する。
\item 積と整合する実際の二周期同型 $\beta:D\xrightarrow{\sim}D(2)$ がある。
\end{enumerate}
これらは補題~\ref{lem:relative-geometric-cubic-flatness} の相対 cubic 解消から得られる入力である。
最後の条件は、幾何ファイバーごとの二周期同型の存在より強い。
本節の積の規約は
\[
 B_n=D_{0,2n},\qquad a\star b=a\,\beta^m(b)
 \quad(a\in B_m,\ b\in B_n)
 \tag{UH-2}\label{hull-eq:product}
\]
である。右軌道積 $a*b=s^n(a)\circ b$ と比較する際は、
式 (\ref{eq:cubic-tails-right-convention}) の
$C_{ij}=D_{-j,-i}$、$\beta_C=s^{-1}$ を用いて規約を合わせる。
偶数環と右軌道環の積の一致は定理~\ref{thm:quadric-finite-type-effectivity} の証明と同じである。

$U_i=V_i\otimes V_{i+1}$ および
\[
 S_i=R_i\otimes V_{i+3}+V_i\otimes R_{i+1}
       \subset U_i\otimes U_{i+2}
 \tag{UH-3}\label{hull-eq:relations}
\]
と置く。$S_i$ は階数七の直和因子である。
$U=U_0$ とし、$U\otimes U\xrightarrow{1\otimes\beta}U_0\otimes U_2$
に沿って $S_0$ を $U^{\otimes2}$ の部分加群とも考える。
この規約で $B=T_R(U)/(S_0)$ である。

\begin{theorem}[一般二周期 cubic 族の canonical AS4 被覆]
\label{thm:universal-cubic-hull}
上の仮定の下で、基底変更に自然な全射
\[
 \Phi_i:S_i\longrightarrow W_i,\qquad
 \Phi_i|_{W_i}=2\,\mathrm{id}_{W_i}
 \tag{UH-4}\label{hull-eq:Phi}
\]
があり、$H_i=\ker\Phi_i$ は階数六の直和因子である。
$\beta(H_i)=H_{i+2}$ であり、
\[
 A=T_R(U)/(H_0)\longrightarrow B
 \tag{UH-5}\label{hull-eq:hull}
\]
は次を満たす。
\begin{enumerate}
\item 各 $A_n$ は階数 $\binom{n+3}{3}$ の有限 projective 加群であり、
  任意の基底変更後にも同じ提示と次数別比較が成立する。
\item $P=A^!$ は階数列 $(1,4,6,4,1)$ の有限 projective な二次代数で、
  $P_n=0$ $(n\ge5)$ である。$P$ は相対 Koszul Frobenius であり、
  $A$ の左右および双加群 Koszul 解消は長さ四で、任意の基底変更後にも完全である。
\item $G=S_0/H_0\cong W_0$ は $A_2$ の直和因子である。
  $J=\ker(A\to B)$ に対し、全次数で普遍的な同型
  \[
   G\otimes_R A_{n-2}\xrightarrow{\sim}J_n
     \xleftarrow{\sim}A_{n-2}\otimes_R G
   \tag{UH-6}\label{hull-eq:regular-normal}
  \]
  がある。従って $A/(G)=B$、$GA=AG$ であり、$G$ の任意の局所生成元は
  左右正則正規な二次の元である。
\item 任意の可換 Noetherian $R$-代数 $R'$ に対して、$A_{R'}$ は左右 Noetherian である。
\end{enumerate}
各幾何ファイバーの被覆は Van den Bergh の canonical hull を
指定された二周期に沿って graded 化したものに一致する。
ここで Serre 平方根に関する仮定は用いない。
\end{theorem}

Noetherian 性には定理~\ref{thm:relative-geometric-cubic-noetherian} を用いる。
それ以外の結論は本節の有限 projective と tensor の議論による。
相対 AS4 という語は基底の解消を指し、特異な $R$ 上で環 $A$ 自身の
絶対大域次元が四であると主張しない。

\subsection{逆 coevaluation と次数三の正規性}

\subsubsection{評価の規約と線束値の写像}

\eqref{hull-eq:coevaluation} の逆を使うと
\[
 R_i\otimes V_{i+3}\cong
 W_i\otimes V_{i+3}^{\vee}\otimes V_{i+3}\longrightarrow W_i
\]
および
\[
 V_i\otimes R_{i+1}\cong
 V_i\otimes W_i\otimes V_i^{\vee}\longrightarrow W_i
\]
という二つの評価写像が得られる。
両者の $W_i$ への制限は階数二の恒等写像の trace、すなわち $2\,\mathrm{id}$ である。
従って交わりで一致し、\eqref{hull-eq:Phi} を定める。
この全射は $W_i\hookrightarrow S_i$ の半分を切断に持ち、
$S_i=H_i\oplus W_i$ が成り立つ。
構成は有限 projective の評価と普遍的な交わりのみを用いるため、
任意の基底変更および $D$ の同型に自然である。

局所生成元 $w_i$ を選び $\Phi_i(x)=\phi_i(x)w_i$ と書くと、
\[
 w_i=\sum_{a=1}^2r_a\otimes v_a
 \quad\Longrightarrow\quad
 \phi_i(r_a\otimes v_b)=\delta_{ab}.
 \tag{UH-7}\label{hull-eq:inverse-coeval}
\]
他方の分解でも同じ規約を用いる。
特に $\phi_i(w_i)=2$ であり、$w_i$ を $u w_i$ に変えると
$\phi_i$ は $u^{-1}\phi_i$ になる。
従って本来の大域対象は単一の scalar-valued functional ではなく、
\eqref{hull-eq:Phi} とその双対線 $W_i^{\vee}\subset S_i^{\vee}$ である。

\subsubsection{四周期と二周期の区別}

$W_0,W_1$ を局所的に自明化し、$w_{i+2}=\beta(w_i)$ と選ぶ。
$w_i\in V_i\otimes R_{i+1}$ と
$w_{i+1}\in R_{i+1}\otimes V_{i+4}$ の二つの非退化 tensor の比較により、
一意な同型 $\theta_i:V_i\xrightarrow{\sim}V_{i+4}$ が定まる。
具体的には $w_i=\sum_a v_a\otimes r_a$ と書いたとき
$w_{i+1}=\sum_a r_a\otimes\theta_i(v_a)$ とする。
四回の巡回から $\theta_i\otimes\theta_{i+1}\otimes\theta_{i+2}$ は
$R_i$ を $R_{i+4}$ に移す。
従って $\theta:D\to D(4)$ は代数同型である。
自然性から
\[
 \theta_{i+2}\beta_i=\beta_{i+4}\theta_i
 \tag{UH-8}\label{hull-eq:commute}
\]
であり、選んだ生成元の下で $\theta$ と $\beta$ は可換する。
この $\theta$ は \cite[\S5.2, (5.2)]{VdB2011} の canonical 四周期である。
一般に $\theta=\beta^2$ は成立しない。

\begin{lemma}[偶数環の相対 Koszul 性]
\label{hull-lem:even-koszul}
$B$ は相対 Koszul であり、$E=B^!$ の全次数は有限 projective で、
任意の基底変更と両立する。階数列は
\[
 (\rank E_n)_{n\ge0}=(1,4,7,8,8,\ldots).
 \tag{UH-9}\label{hull-eq:E-ranks}
\]
\end{lemma}
\begin{proof}
体上で、標準 cubic 解消を偶数頂点と奇数頂点に分ける。
奇数頂点の加群には、$V$ と $R$ を交互に用いる偶数環上の線型解消が得られる。
偶数部分の二つの短完全列に Tor を適用すると、偶数環の単純加群の
Tor は対角次数に集中する。これは \cite[Proposition 2.1]{BauwensVanDenBergh1998}
の証明に各 $V_i,R_i,W_i$ の添字を付けたものであり、実際の一周期を必要としない。
従って全幾何ファイバーの $B$ は Koszul である。

$B_1,B_2,B_3$ は仮定により有限 projective である。
\cite[Chapter 6, Theorem 7.3, Propositions 7.1--7.2]{PolishchukPositselski2005}
を各局所環で適用すれば、関係部分加群の全交わりと商が直和因子となる
（補題~\ref{lem:relative-koszul-degree-three}）。
これにより $E$ の次数別 projectivity、任意の基底変更との整合、
左右の普遍 Koszul 解消が得られる。
各ファイバーでは $\dim B_n=(n+1)^2$ だから
$h_E(t)=h_B(-t)^{-1}=(1+t)^3/(1-t)$ となり、\eqref{hull-eq:E-ranks} が従う。
\end{proof}

\begin{lemma}[普遍的な次数三分解]
\label{hull-lem:K3}
$E_3^{\vee}$ の対応する添字付き成分は
\begin{align*}
 K_{3,i}
 &= (S_i\otimes V_{i+4}\otimes V_{i+5})
       \cap(V_i\otimes V_{i+1}\otimes S_{i+2})\\
 &= (R_i\otimes R_{i+3})
       \oplus(V_i\otimes W_{i+1}\otimes V_{i+5}).
 \tag{UH-10}\label{hull-eq:K3}
\end{align*}
この同定は任意の基底変更と両立する。
\end{lemma}
\begin{proof}
右辺の二つの部分加群は $W_{i+1}$ の二つの包含により左辺に入る。
左辺は補題~\ref{hull-lem:even-koszul} により階数八の projective 加群であり、
右辺は階数四の projective 加群二つの直和からの射である。
従って各幾何ファイバーで直和であることを示せば十分である。

体上で二つの像の交わりを取ると、
$R_iV_{i+3}\cap V_iR_{i+1}=W_i$ により、それは
$W_iV_{i+4}V_{i+5}$ に入る。
一方、右 cubic 解消の次数六の部分は
\[
 W_i\otimes V_{i+4}\otimes V_{i+5}
 \hookrightarrow R_i\otimes D_{i+3,i+6}
\]
を単射にする。
右辺は $R_iV_{i+3}V_{i+4}V_{i+5}$ を $R_iR_{i+3}$ で割ったものだから、
$W_iV_{i+4}V_{i+5}\cap R_iR_{i+3}=0$ である。
従って交わりは零となる。次元八より同型である。
この議論は \cite[(2.7), Lemmas 2.3--2.4]{BauwensVanDenBergh1998} の添字付きの証明でもある。
有限 projective の間の同型判定を各ファイバーで行えば主張を得る。
\end{proof}

\begin{lemma}[functional 線の普遍正規性]
\label{hull-lem:normal}
\eqref{hull-eq:Phi} の双対は可逆加群 $N=W_0^{\vee}\subset E_2$ を与え、
$NE=EN$ が任意の基底変更後に成立する。
\end{lemma}
\begin{proof}
局所的な $w_i$ の選択の下で
\[
 c_L=\phi_i\otimes1:K_{3,i}\longrightarrow U_{i+4},\qquad
 c_R=1\otimes\phi_{i+2}:K_{3,i}\longrightarrow U_i
\]
と置く。$\Theta_i=\theta_i\otimes\theta_{i+1}$ とすると
\[
 c_L=\Theta_i c_R
 \tag{UH-11}\label{hull-eq:normal-contraction}
\]
である。これはファイバーでのみ成立する式ではなく、次の tensor 計算による恒等式である。

$\theta$ により四つ先の空間を同一視する。
$V_iW_{i+1}V_{i+5}$ 上では二つの縮約はいずれも
coevaluation と evaluation の合成なので同じ恒等射になる。
第一成分では、四つの空間の基底を $u_a,v_b,x_c,y_d$ とし、
$w_i=\sum t_{abcd}u_a v_b x_c y_d$ と書く。
\[
 r_d=\sum_{a,b,c}t_{abcd}u_av_bx_c\in R_i,\qquad
 t_c=\sum_{d,a,b}t_{abcd}y_du_av_b\in R_{i+3}
\]
はそれぞれ基底である。両縮約は $r_d\otimes t_c$ を
$\sum_{a,b}t_{abcd}u_av_b$ へ送るため、\eqref{hull-eq:normal-contraction} が成立する。

\eqref{hull-eq:product} に沿って graded 化する。
$T=\beta^{-2}\theta$ が $U$ に誘導する可逆写像を同じ記号で表すと、
\eqref{hull-eq:normal-contraction} の双対は
\[
 \phi\xi=(\xi\circ T)\phi\qquad(\xi\in E_1)
 \tag{UH-12}\label{hull-eq:graded-normality}
\]
となる。$T$ は代数自己同型から来るので、その双対は $E$ の自己同型を定める。
一次生成性から $\phi E=E\phi$ を得る。
生成元を替えると $\phi$ が単元倍されるだけなので、これは大域的な線 $N$ の正規性を意味する。
計算は基底上の恒等式であり、任意の基底変更後も同じである。
\end{proof}

\subsection{全幾何ファイバーでの canonical hull との同定}

本節のみ基礎環を代数閉体 $K$ とする。
\cite[Theorem 6.1]{VdB2011} により、canonical 四周期 $\theta$ の逆の
正規化自己同型の類を持つ、偶数 $\mathbb Z$-代数 $D^{\mathrm{ev}}$ の hull が存在する。
以下では、そこに現れる関係空間が実際に $H_i=\ker\phi_i$ であることを証明する。
これは、既知の hull の存在だけから候補の同定を省略しないための一段である。

\begin{lemma}[正規化式が決める評価の一意性]
\label{hull-lem:fibre-identification}
canonical hull の六次元関係空間を $H'_i\subset S_i$ とする。
このとき $H'_i=H_i$ である。
\end{lemma}
\begin{proof}
$H'_i=\ker\psi_i$ となる非零写像 $\psi_i:S_i\to K$ を選ぶ。
$g_i\in S_i/H'_i$ を $\psi_i(g_i)=1$ で正規化する。
この記法では、$S_i$ の元の hull における像は $\psi_i$ 倍の $g_i$ である。
canonical normalizer の意味から、ある $\lambda_i\in K^\times$ に対し
\[
 g_i u=\lambda_i\Theta_i^{-1}(u)g_{i+2}
 \qquad(u\in U_{i+4})
 \tag{UH-13}\label{hull-eq:hull-normalizer}
\]
である。scalar 自己同型の不定性は $\lambda_i$ に吸収した。

$z\in K_{3,i}$ を二通りに hull の次数三へ送れば
\[
 g_i(\psi_i\otimes1)(z)
   =(1\otimes\psi_{i+2})(z)g_{i+2}
\]
となる。$g_{i+2}$ の右正則性と \eqref{hull-eq:hull-normalizer} より
\[
 (\psi_i\otimes1)(z)
   =\lambda_i^{-1}\Theta_i(1\otimes\psi_{i+2})(z).
 \tag{UH-14}\label{hull-eq:psi-normalizer}
\]

これを第二直和成分 $V_iW_{i+1}V_{i+5}$ に制限する。
coevaluation によって、左辺は $F_i\otimes\mathrm{id}$、右辺は
$\lambda_i^{-1}\mathrm{id}\otimes G_i$ という形に書ける。
ここでは $\theta_i,\theta_{i+1}$ により始点と終点を同一視した。
二つの非零有限次元空間の endomorphism が
$F\otimes\mathrm{id}=\mathrm{id}\otimes G$ を満たせば、
行列単位との比較により $F,G$ は同じ scalar の恒等写像である。
従って
\[
 \psi_i|_{V_iR_{i+1}}=a_i\phi_i|_{V_iR_{i+1}},\qquad
 \psi_{i+2}|_{R_{i+2}V_{i+5}}
       =b_i\phi_{i+2}|_{R_{i+2}V_{i+5}}
\]
となる。$i$ を二つ戻した式も用いれば、$\psi_i$ の二つの制限は
それぞれ逆 coevaluation の scalar 倍である。
交わり $W_i$ 上で両者は同じ $\psi_i(w_i)$ を与え、
$\phi_i(w_i)=2\ne0$ なので二つの scalar は一致する。
よって $\psi_i=a_i\phi_i$ である。
$\psi_i\ne0$ から $a_i\ne0$、従って $H'_i=H_i$ を得る。
\end{proof}

\begin{remark}[Bauwens--Van den Bergh の式との直接照合]
\cite[Proposition 2.7, (2.11), Remark 3.3]{BauwensVanDenBergh1998} は graded cubic の場合、
\[
 w=x^tf=(Qf)^tx,\qquad
 \psi(f_i x_j)=\eta(\Lambda^{-1})_{ji},\qquad
 \psi(x_i f_j)=\eta(\Lambda Q^{-t})_{ij}
\]
と書く。canonical choice は $\Lambda=Q^t$ であるから
\[
 \psi(x_i f_j)=\eta\delta_{ij},\qquad
 \psi(f_i x_j)=\eta(Q^{-1})_{ij}.
 \tag{UH-15}\label{hull-eq:BV-convention}
\]
これは $w$ の二つの分解に対する \eqref{hull-eq:inverse-coeval} と正確に一致する。
原典では偶数環の元の次数を元の cubic 環の次数で記している箇所があるため、
そこに現れる次数四の functional は本節の $E_2$ である。
一般 $\mathbb Z$-代数での同定は補題~\ref{hull-lem:fibre-identification} が直接与える。
\end{remark}

\begin{corollary}[体上の双対と正則性]
\label{hull-cor:fibre-dual}
任意の実際の二周期 $\beta$ に対して、$T(U)/(H_0)$ は canonical hull の
graded 化である。$\phi\in E_2$ は左右正則正規で、
$E/(\phi)$ は階数列 $(1,4,6,4,1)$ の Koszul Frobenius 代数である。
\end{corollary}
\begin{proof}
$H_i$ の構成の自然性により $\beta(H_i)=H_{i+2}$ であるから、
二周期は hull の提示そのものに持ち上がる。
これは \cite[Lemmas 6.4--6.5, Proposition 6.6]{VdB2011} の構成と一致する。
canonical hull は標準四次元 AS 正則なので Koszul であり、
その二次双対 $P$ は Frobenius、Hilbert 級数 $(1+t)^4$ を持つ。
二次提示から $P=E/(\phi)$ である。
補題~\ref{hull-lem:normal} より $\phi$ は正規である。
\[
 h_E(t)=\frac{(1+t)^4}{1-t^2},\qquad h_P(t)=(1+t)^4
\]
なので、各次数の商と積の次元を比べると
$E_{n-2}\xrightarrow{\phi\cdot}E_n$ および
$E_{n-2}\xrightarrow{\cdot\phi}E_n$ は単射となる。
これで正則性が従う。原典でも \cite[Propositions 2.7, 3.2]{BauwensVanDenBergh1998} が同じ双対関係を示す。
\end{proof}

\subsection{非被約基底上での相対化}

再び一般の Noetherian $R$ に戻る。
定理~\ref{thm:universal-cubic-hull} の証明を完結する。

\subsubsection{まず双対の商を平坦にする}

補題~\ref{hull-lem:even-koszul} より $E_n$ は有限 projective である。
補題~\ref{hull-lem:normal} の正規線 $N=W_0^{\vee}$ による左右乗法
\[
 N\otimes E_{n-2}\longrightarrow E_n,
 \qquad E_{n-2}\otimes N\longrightarrow E_n
 \tag{UH-16}\label{hull-eq:E-inject}
\]
は、系~\ref{hull-cor:fibre-dual} により全幾何ファイバーで単射である。
有限 projective 加群間の写像のファイバー単射性から、
これらは局所直和因子への単射になる。
従って
\[
 P=E/(N)
 \tag{UH-17}\label{hull-eq:P}
\]
の全次数は有限 projective であり、商は任意の基底変更と可換する。
\eqref{hull-eq:E-ranks} の差から階数列は $(1,4,6,4,1,0,\ldots)$ となる。
ここでは商のファイバー次元が一定であることだけから平坦性を推論していない。
既に projective である加群間の写像 \eqref{hull-eq:E-inject} を用いている。

$P$ は二次代数であり、$P_1,P_2,P_3$ は projective、全幾何ファイバーは Koszul である。
従って再び \cite[Chapter 6, Theorem 7.3]{PolishchukPositselski2005} を \emph{$P$ に}適用する。
全ての関係交わりが直和因子となるため、二次双対
\[
 P^!=T_R(U)/(H_0)=A
\]
の全次数も有限 projective で、任意の基底変更に整合する。
全ファイバーでの canonical hull の Hilbert 級数により
$\rank A_n=\binom{n+3}{3}$ となる。
この順序により、$A_3$ の局所自由性を証明前に仮定する循環を避けている。

\subsubsection{Frobenius、AS 解消、二次切断}

$L=P_4$ は可逆加群である。
各積の adjoint
\[
 P_i\longrightarrow\operatorname{Hom}_R(P_{4-i},L)
 \qquad(0\le i\le4)
 \tag{UH-18}\label{hull-eq:Frobenius}
\]
は有限 projective 加群間の射で、全幾何ファイバーで同型である。
従って同型であり、任意の基底変更後もそうである。
これが相対 Frobenius 性である。
Koszul 複体の各内部次数は有限な projective 複体で、
そのファイバーは体上の Koszul 解消だから分裂完全となる。
$P_i=0$ $(i>4)$ より、左右および対角双加群の解消は長さ四である。
Frobenius pairing による双対化から
\[
 \underline{\operatorname{Ext}}_A^j(R,A)
 \cong
 \begin{cases}
 L\otimes_RR(4),&j=4,\\
 0,&j\ne4
 \end{cases}
 \tag{UH-19}\label{hull-eq:AS}
\]
を左右に得る。ここで表示は対応する側の加群としてのものとし、
双加群同定には通常の Frobenius 自己同型を含める。
解消は任意の基底変更後にも完全なので、これらは相対 AS4 条件である。

$A\to B$ は提示から全射である。
$A_n,B_n$ が有限 projective だから $J_n=\ker(A_n\to B_n)$ も projective であり、
核は任意の基底変更と可換する。
$G=S_0/H_0\cong W_0$ は $J_2$ と同一視される。
左右乗法 \eqref{hull-eq:regular-normal} は有限 projective 加群間の射で、
各幾何ファイバーでは canonical hull の正則二次切断の同型である。
従って全次数で同型であり、任意の基底変更後も同型である。
これにより正則性、正規性、商 $A/(G)=B$ を同時に得る。
特に局所生成元 $g$ について $gA=Ag$ と両側正則性があるので、
$ga=\nu(a)g$ によって定まる写像 $\nu$ は全単射の代数自己同型である。
規約 \eqref{hull-eq:graded-normality} では、$\nu|_U=T^{-1}$ となる。
実際、$\phi(c)=1$ となる $c\in S_0$ で $g$ を表し、
$c\otimes u-\nu(u)\otimes c=h_L+h_R$、
$h_L\in H_0\otimes U$、$h_R\in U\otimes H_0$ と書く。
すると $z=c\otimes u-h_L=\nu(u)\otimes c+h_R\in K_3$ であり、
二つの縮約はそれぞれ $u$ と $\nu(u)$ になる。
\eqref{hull-eq:normal-contraction} を graded 化した等式から $u=T\nu(u)$ を得る。
この確認も幾何ファイバーへの還元を用いない。

\subsubsection{有限判定と Noetherian 性}

定理~\ref{thm:relative-finite-as-cover} の五条件は、本構成では全て基底上で成立する。
具体的には
\[
 \rank A_3=20,\quad\rank A_4=35,\qquad
 \rank P_3=4,\quad\rank P_4=1,\quad P_5=0,
\]
左右の $G$ 乗法の像の一致、次数一・二に対する局所直和因子への単射、
および \eqref{hull-eq:Frobenius} である。
従って canonical な階数六部分束 $H_0\subset S_0$ は、
被覆の有限判定部分スキームへの実際の切断を与える。
全ファイバーで有限条件が成り立つという主張のみで代用していない。

最後に任意の Noetherian $R'$ に移る。
系~\ref{cor:relative-periodic-even-noetherian} より $B_{R'}$ は左右 Noetherian である。
$G_{R'}$ を有限アフィン開被覆上で自明化する。
補題~\ref{noeth-lem:normal-lift} の議論を適用する。
正則正規な正次数元 $g$ に対して、$\operatorname{gr}_{(g)}A_{R'}$ は
$B_{R'}$ 上の skew polynomial 環となるため左右 Noetherian である。
固定内部次数では $(g)$ 進フィルトレーションは有限なので、
初期斉次生成元を有限個持ち上げる還元は有限回で止まる。
これにより斉次イデアルは有限生成、最高内部次数による還元で一般のイデアルも有限生成となる。
基底の有限開被覆から忠実平坦に降下すれば $A_{R'}$ は左右 Noetherian である。
これで定理~\ref{thm:universal-cubic-hull} の全ての主張を得た。

\subsection{Serre 平方根の条件は独立である}

\begin{remark}[二周期と反 Serre 関手]
任意の実際の二周期 $\beta$ が上の AS4 被覆を持つことと、
そこから来る exact 自己同値 $s$ が
\[
 s^2\simeq S^{-1}[2]
 \tag{UH-20}\label{hull-eq:Serre-root}
\]
を満たすことは別である。
後者は本定理の仮定にも結論にも含まれない。
有限 foundation $\Lambda$ 上で $s$ が可逆 perfect 双加群 $P_s$ により表され、
$\mathcal S_\Lambda=\mathbf R\operatorname{Hom}_R(\Lambda,R)$ が Serre 双加群である場合、
正確な追加条件は
\[
 P_s\mathbin{\otimes^{\mathbf L}}_\Lambda P_s
    \simeq\mathcal S_\Lambda^{-1}[2]
 \quad\text{in }\operatorname{Perf}(\Lambda^e).
 \tag{UH-21}\label{hull-eq:Morita-Serre}
\]
右辺の逆は可逆双加群の Morita 逆である。
根・根同型を fppf 局所の存在条件として扱う定義では、
基底線束 $\ell$ を許した
$P_s\mathbin{\otimes^{\mathbf L}} P_s\simeq\mathcal S_\Lambda^{-1}[2]\otimes_R\ell$ を
局所自明化して同じ条件を表せる。
この比較は定理~\ref{thm:bimodule-line-comparison} と
系~\ref{cor:foundation-antiserre-comparison-locus} で認識する対象である。
単に対象ごとに $s(F_0)\cong F_2$, $s(F_1)\cong F_3$ であることだけでは
\eqref{hull-eq:Morita-Serre} は従わない。
\end{remark}

\begin{example}[period two は自動的な Serre 根ではない]
$Q=\mathbb P^1\times\mathbb P^1$、$\sigma\in\operatorname{Aut}(Q)$ を
$\sigma^2\ne1$ かつ両 ruling を保つ自己同型とする。
例えば標数零で
$\sigma([x_0:x_1],[y_0:y_1])=([2x_0:x_1],[y_0:y_1])$ とできる。
$s=(-\otimes\OO_Q(1,1))\circ\sigma^*$ は標準 cubic helix の各対象を二項先へ送り、
実際の二周期を与える。
$s^{-n}\OO_Q\cong\OO_Q(-n,-n)$ なので $(\OO_Q,s)$ は AZ ample でもある。
従ってその偶数軌道環には本定理の canonical AS4 被覆がある。
しかし
\[
 s^2=(-\otimes\OO_Q(2,2))\circ\sigma^{2*},\qquad
 S^{-1}[2]=(-\otimes\OO_Q(2,2))
\]
であり、動かされる点の skyscraper sheaf を比べれば同型でない。
従って同じ指定二周期を H2 の witness とするには、
\eqref{hull-eq:Serre-root} を別途確かめる必要がある。
この例は underlying $(\coh Q,\OO_Q)$ が H2 に属することを否定しない。
\end{example}

\paragraph{一般 cubic 点への適用で残る入力.}
本節により、実際の geometric cubic 二周期族が構成された後の AS4 被覆は
Q 固有の障害ではなくなる。
一般の内在的 H2 族に適用するには、構造対象を含む適合する四項 foundation、
その上で代数化された実際の Morita 根、
および \eqref{hull-eq:Morita-Serre} と全軌道環の比較を供給する必要がある。
これらの存在範囲および根を忘れた相対実現部分層の代数性は、本節では主張しない。

